\section{Implementation}

\subsection{The transforms in detail}

\textbf{Canonicalize} reorders the operands of commutative instructions into a
stable form: constants move to the right, and other operands are ordered by their
position in the function. Swapping the operands of a commutative operation never
changes its result, so the rewrite is unconditionally safe, and it makes both
later passes more effective at once.

\textbf{Algebraic simplify} folds integer identities (an add of zero, a subtract
of self, a multiply by one or zero, the and, or, and xor identities, and a shift
by zero) and strength reduces a multiply by a constant power of two into a shift,
carrying the no signed wrap and no unsigned wrap flags onto the shift. It works
only on scalar integers. Floating point is left alone because folding a floating
point identity is unsound without fast math flags, and the predicates only match
integer constants, so the pass never even considers it. Vectors, volatile, and
atomic operations are likewise untouched.

\textbf{Local common subexpression elimination} walks each block once and
replaces a pure computation that repeats an earlier identical one with a
reference to the first. Two instructions match by a structural key that orders
commutative operands, which is why canonicalization runs first. Loads are never
deduplicated, because two identical loads can observe different memory.

\textbf{Dead code elimination} removes trivially dead instructions with a
worklist, so that deleting one instruction can make its operands dead in turn.
Anything with a side effect (a store, a volatile or atomic access, a call not
known to be free of effects) is kept.

\subsection{Correctness by construction and by test}

Every pattern has a lit test that pins the exact IR before and after, and every
transform has negative tests: a floating point add of zero, a multiply by a non
power of two, a volatile load, a store, and a call are all shown to survive
untouched. Beyond the pattern tests, an invalidation test caches an analysis
result, runs a transform through a real pass manager, and asserts the result
recomputes rather than serving the stale cache. This test fails if a transform
over claims what it preserves, which is the single most common way to get the new
pass manager contract wrong.
